% ---------------------------------------------------------------------
\section{Compared methods, training and statistics}\label{sec:training-supp}
% ---------------------------------------------------------------------

\subsection{Implementation of the compared methods}\label{app:baselines}

Every method and every evaluation pass receives the shared admitted requirement of \Sec{sec:pipeline}, and shortfall is always measured against it.
Methods without the safety layer execute only the physical clip $\clip(x_i,\pmn_i,u_i)$, and variants with the layer pass their output as the raw action $x$ through \Alg{alg:safety}, so they inherit Theorem~\ref{M-thm:invariance}.

\textbf{Penalty-based methods.}
Attention+penalty, MLP+penalty and the variant without departure bound learn from the reward $r(t)$ minus two penalties.
They are $w$ times the departure shortfall of the EVs that leave at step $t$ and $w$ times the transformer overload energy of the step.
Shortfall is measured against the admitted requirement, overload uses the untightened limits and the reported powers, and the weight $w\in\{1,10,100\}$\,\texteuro/kWh is tuned.
At the evaluation weight of 10\,\texteuro/kWh, more than eleven times the highest hourly price, no violation pays at the optimum of the penalised objective, although learned policies can still violate.

Attention+penalty and MLP+penalty use $\beta=0$, and their actor loss masks only tokens whose action cannot influence the commanded power.
The variant without departure bound runs the safety layer with $l_i=\pmn_i$, so it keeps $\beta=0.01$ and masks fallback steps like \method{}.

\textbf{Attention+PPO-Lagrangian.}
This baseline solves the constrained problem $\max J_r$ subject to expected episode shortfall and overload energy of at most zero \cite{altman2021constrained,tessler2018reward}.
The attention policy runs with the physical clip and the profit reward, and the critic has three heads, one for the value and one for each cost, each with its own advantage estimate.
PPO uses the combined advantage $(A_r-\Lambda_EA_E-\Lambda_PA_P)/(1+\Lambda_E+\Lambda_P)$, normalised per minibatch.
Once per rollout, $\Lambda_k\gets\clip(\Lambda_k+\eta_\Lambda(\hat J_k-0),0,1000)$ for $k\in\{E,P\}$, with $\Lambda_k=0$ at the start and the step size $\eta_\Lambda\in\{0.5,2,8\}$ tuned.
The cost limit of zero matches the hard requirements of the problem, and the cap of 1000 keeps the multipliers finite while violations persist, and the training logs record their trajectories.

\textbf{Size-agnostic and fixed-size networks.}
Deep Sets+safety layer embeds each charger token with a shared network and adds the masked mean of all embeddings and the global token as context.
It decodes one Gaussian per token, so no token attends to another.
The MLP baselines map the zero-padded concatenation of 20 charger tokens and the global token through two hidden layers of 640 units.
Their width matches the parameter count of the attention actor (614,440 versus 610,370), and the MLP baselines are evaluated and reported only at $N\le20$.


\textbf{Rules.}
The rules as fast as possible (AFAP), as late as possible (ALAP), round robin and LLF are price-unaware and charge only.
For each EV, sLLF \cite{chen2022smoothed} computes the next-step laxity $\chi_i^0=(E_i-\Bh_i(\tau_i-1))/(\Pnom_i\Delta t)$ and the rate $p_i=\clip(\Pnom_i(L-\chi_i^0),0,u_i)$, with a common level $L$ found by bisection so that $\sum_ip_i=\max(0,\Ctil)$.
It is the published rule with the peak rate $\Pnom_i$ and a cap that includes the taper.
The price-threshold rule compares the current price with the median of all prices seen so far together with the known price window.
Below the median it allocates $\Ctil(t)$ by LLF, above the median it requests the maximum V2G discharge, and otherwise it idles.

\textbf{Model-predictive baselines.}
eMPC-HiGHS is our re-implementation of the eMPC of EV2Gym \cite{diazlondono2026open} as a receding-horizon net-power LP with HiGHS.
Its LP uses per-EV limits, exact first-step bounds including the taper and the admitted requirements with an elastic departure constraint of 1000\,\texteuro/kWh.
Prices and arrivals are known exactly within its horizon of 10 steps, whereas load and PV come from EV2Gym's noisy forecasts.
A 24-h horizon variant and the original Gurobi-based eMPC were planned but not run.

\textbf{External learned baselines.}
The heterogeneous graph policy (EV-GNN) \cite{orfanoudakis2025scalable} was planned as an external size-agnostic baseline, to be run from its public code at commit \texttt{01aa1b7a} with its own training script, but this study did not include it.
In addition, the graph-enhanced decision transformer (GNN-DT) \cite{orfanoudakis2026graph} has public code under an open-source licence (DT4EVs, commit \texttt{8b327d12}).
GNN-DT trains offline on trajectories that an optimiser generates for each scenario inside a modified simulator copy, and porting this pipeline to our scenarios was not done.

\subsection{Training loop}\label{app:ppo}

\Alg{alg:ppo} lists the training loop of \method{}, in which the safety layer of \Alg{alg:safety} acts as part of the environment.
Each environment writes into its own buffer, and at update time all token tensors are right-padded to the largest $N$ in the batch with masked entries.
Padding changes neither the actor outputs of real tokens nor the value (Lemma~\ref{lem:equivariance}), and the log-probability, entropy and regulariser sums run over masks only.
The learning reward is the profit divided by $N r^{\sm{ref}}$, so that returns of stations with different $N$ have the same scale, and every reported result uses the undivided profit.

\begin{algorithm}[t]
\caption{PPO training of \method{} with a variable number of chargers}\label{alg:ppo}
\begin{algorithmic}[1]
\Require 16 environments (8 with $N=10$, 8 with $N=20$); actor $\pi_\theta$ and critic $V_\omega$
\While{fewer than $3\times10^{6}$ environment steps}
  \For{128 steps in every environment}
    \State observe tokens, $l$, $u$, $\Ctil$, $\Crev$ and the feasibility flag
    \State sample $a\sim\pi_\theta(\cdot\mid s)$; $x\gets P^{\sm{ref}}a$
    \State $p\gets\Pi_t(x)$ if feasible, else fallback rule \Comment{\Alg{alg:safety}}
    \State step the simulator; $r\gets$ realised profit of the step (\Eq{M-eq:objective}); store $r/(N r^{\sm{ref}})$, $a$, $\log\pi_\theta(a\mid s)$ over tokens with $u_i-l_i>10^{-9}$
  \EndFor
  \State compute advantages with GAE ($\gamma=0.99$, $\lambda_{\sm{GAE}}=0.95$); pad tokens to the largest $N$ in the batch
  \For{10 epochs, minibatches of $n_{\sm{mb}}$ transitions}
    \State clipped surrogate ($\epsilon_{\sm{clip}}=0.2$) on feasible transitions with normalised advantages, divided by $n_{\sm{mb}}$
    \State loss $\gets-$surrogate $+\,c_{\sm{vf}}\cdot$value loss $-\,c_{\sm{ent}}\cdot$entropy $+\,\mathcal L_{\sm{reg}}$ (\Eq{M-eq:reg}, feasible steps only)
    \State Adam step with gradient-norm clipping 0.5; stop the epochs early if the approximate KL divergence exceeds 0.03
  \EndFor
\EndWhile
\end{algorithmic}
\end{algorithm}

\subsection{Hyperparameters and tuning}\label{app:hyper}

\Tab{tab:hyper} lists the hyperparameters.
Each of the seven tuned learned methods (\method{}, Attention+penalty, Attention+PPO-Lagrangian, Attention+safety layer without departure bound, Deep Sets+safety layer, MLP+safety layer and MLP+penalty) received 12 configurations.
They consist of the default point and 11 points sampled without replacement from the grid of \Tab{tab:hyper} with a fixed sampling seed per method.
Each configuration is trained on the training mixture for $0.5\times10^{6}$ steps with seed 0 and evaluated on 24 validation episodes that cycle through the training configurations, with seeds disjoint from the test replays.
Selection uses the mean validation $J^{\sm{pen}}$ at $w=10$\,\texteuro/kWh for every method.
The conformal variant and the reference trained at $N=50$ inherit the selection of \method{}, and $\beta$ and the architecture are fixed for all methods with the safety layer.
Tuning ran 84 training and 84 validation jobs, and the selected configurations are released with the code.

EV2Gym could not build some training episodes, because negative residential loads gave a negative scale of the forecast noise.
The training seeds were therefore redrawn, once per tuning job and twice per training job, and the test replays were built unchanged.

\begin{table}[t]
\centering
\caption{Hyperparameters and tuning grid.}
\label{tab:hyper}
\footnotesize
\setlength{\tabcolsep}{3pt}
\begin{tabular}{@{}l>{\raggedright}p{6.2cm}>{\raggedright\arraybackslash}p{6.4cm}@{}}
\toprule
Group & Parameter & Value \\
\midrule
Network & width $d_m$, layers $n_L$, heads $n_H$ & 128, 3 (pre-LayerNorm), 4 \\
 & feed-forward width, dropout & 512 (GELU), 0 \\
 & actor head (shared), critic readout & $128\to64\to2$; attention pooling with one query, $128\to64\to1$ \\
 & $\log\sigma$ clamp and initial value & $[-5,0.5]$, $-0.5$ \\
 & parameters (actor / critic) & 610,370 / 676,481; small network 228,067 in total (head width 32) \\
 & precision & float32 networks, float64 safety layer \\
PPO & optimiser & Adam, $\epsilon=10^{-5}$, linear learning-rate decay \\
 & learning rate (tuned) & $\{1,3,10\}\times10^{-4}$ \\
 & minibatch $n_{\sm{mb}}$ (tuned) & $\{256,512,1024\}$ \\
 & entropy coefficient $c_{\sm{ent}}$ (tuned) & $\{0,10^{-3},10^{-2}\}$ \\
 & discount $\gamma$, GAE $\lambda_{\sm{GAE}}$ & 0.99, 0.95 \\
 & clip range $\epsilon_{\sm{clip}}$, value coefficient $c_{\sm{vf}}$ & 0.2, 0.5 \\
 & environments, rollout length & 16 (8 with $N=10$, 8 with $N=20$), 128 \\
 & epochs, target KL, gradient-norm clip & 10, 0.03, 0.5 \\
 & advantage normalisation & per minibatch \\
 & steps per run, seeds & $3\times10^{6}$; 2 (seeds 0 and 1) \\
Penalty reward & training weight $w$ (tuned), $\beta$ & $\{1,10,100\}$\,\texteuro/kWh; $\beta=0$, except 0.01 without departure bound \\
Lagrangian & dual step $\eta_\Lambda$ (tuned), cost limit, $\Lambda_{\max}$ & $\{0.5,2,8\}$, 0, 1000 \\
Safety layer & $P^{\sm{ref}}$, $\beta$ & 11 kW, 0.01 \\
 & $\varepsilon_q$, $\varepsilon_{\sm{sat}}$, tolerance & 0.008 kW, 0.001 kW, $10^{-9}$ kW or kWh \\
Features & $E^{\sm{ref}}$, $c^{\sm{ref}}$, $H_f$ & 100 kWh, 0.25 \texteuro/kWh, 8 steps \\
Evaluation & weights $w_E$, $w_P$ in $J^{\sm{pen}}$ & 10 \texteuro/kWh each, all methods \\
 & test action & deterministic, $a=\mu_\theta$ \\
Conformal & $n_{\sm{cal}}$, $\alpha$ & 5,000 per configuration, $\{0.01,0.05,0.1,0.2\}$ \\
 & quantile predictor & gradient boosting, 200 trees of depth 3, 2,000 training sessions \\
\bottomrule
\end{tabular}
\end{table}

\subsection{Statistical analysis}\label{app:stats}

Each learned method was trained with two seeds, and the reference trained at $N=50$ also with two.
With two seeds, a seed-level $t$-test has one degree of freedom, so no seed-level test or interval is computed.
For each test replay, the value of a learned method is the mean over its two seeds, the seed-averaged policy, and all methods run on the same fixed replays.
Tables report the mean over configurations of the replay means of the seed-averaged policy, followed by the range of the same statistic over the two seeds.
Deterministic rules and eMPC-HiGHS run once and carry no range.
Confidence intervals come from a percentile bootstrap with $10^4$ resamples of the test replays within each configuration, and the interquartile mean (IQM) of the per-episode gap uses $2\times10^3$ resamples \cite{agarwal2021deep}.

Comparisons with \method{} run per station size on the metrics of each results table, and each safety-layer method is also compared with its penalty-based twin.
The primary test is a two-sided Wilcoxon signed-rank test \cite{wilcoxon1945individual} on the replay-level differences of the seed-averaged policies, with the matched-pairs rank-biserial correlation as effect size.
Holm's correction runs over all comparisons of one test file, and each file covers one results table, with up to 351 comparisons per family.
The Wilcoxon test concerns the median shift of the differences, whereas the bootstrap interval concerns their mean.
A few episodes with large penalties can move the mean without moving the median, so the two can disagree, and such comparisons are reported as inconclusive.
Between-seed variability is not captured by any test, and results over few seeds can mislead \cite{henderson2018deep}, so the seed ranges and the seed count are reported as a limitation.

\subsection{Compute}\label{app:compute}

The runs used an AMD Ryzen 7 7700 with 8 cores and 16 threads and an NVIDIA GeForce RTX 5070 Ti with 16 GB.
WSL2 was assigned 12 threads and 31 GB of memory.
As software, Windows 11 Pro with WSL2 Ubuntu 24.04, Python 3.13.15 and PyTorch 2.14.1 with CUDA 13.0 were used.
Preparing scenarios, admission tables and conformal calibrations took 1.6 h, tuning 7.9 h and training of 18 runs 10.2 h.
Evaluation with 298 evaluation and 43 optimum jobs took 2.5 h, and the analysis 6 min.
The core set produced 153,100 episodes of 30 methods, and no job failed.
\Tab{tab:t6} gives the training time per run and the inference time per control step.
Without retraining, the conformal variant is evaluated at all four values of $\alpha$, because Proposition~\ref{M-prop:conformal} holds for any policy.
